% =====================================================================
% QCD confinement vs resonant ternary dynamics (no derivation of 539)
% =====================================================================

\chapter{QCD Confinement versus Resonant Ternary Dynamics}
\label{ch:qcd-vs-resonant}

\section{Confinement in QCD}

Confinement in QCD is the empirical fact that quarks and gluons are never
observed as free asymptotic particles. The asymptotic states of the theory are
colour-singlet hadrons. The interaction between a static quark--antiquark pair
grows linearly with separation at large distance, with a string tension of
order \(1\,\mathrm{GeV/fm}\). When the separation becomes large enough the
string breaks by the creation of a new quark--antiquark pair, producing two
colour-singlet mesons rather than free quarks.

Lattice QCD supplies the strongest theoretical evidence. Wilson loops measured
on large Euclidean lattices obey an area law at intermediate distances; the
extracted string tension scales correctly toward the continuum limit and matches
phenomenological values. Calculations of the static potential, of flux-tube
profiles, and of the spectrum of glueballs and hybrid mesons are all consistent
with confinement. No practical lattice simulation has ever produced free quark
states at zero temperature and low density.

A complete analytic proof of confinement from the continuum QCD Lagrangian is
still lacking. The problem remains one of the Clay Mathematics Institute
Millennium Prize Problems in its Yang--Mills form: the existence of a mass gap
and the related confinement property have not been established rigorously in
four-dimensional continuum non-abelian gauge theory. Many partial results and
dual-superconductivity scenarios exist, but none is regarded as a final proof.

\section{Absence of this apparatus in the resonant map}

\begin{theorem}[No QCD confinement apparatus in \(T_3/T^\sharp\)]
\label{thm:no-qcd-in-map}
None of the following is present in the resonant dynamics of the ternary map
under discussion: a non-abelian gauge connection, Wilson loops, a linear static
potential for colour sources, or a genuine string tension. Modular projections
that keep trajectories on the model's resonant lattice restore residue or charge
conditions; they do not define a gauge field or generate confinement of quarks.
\end{theorem}

\begin{proof}[Proof sketch]
The state space is \(\mathbb{N}\) (or a seeded integer orbit) under \(T_3\) or
\(T^\sharp\). No principal \(G\)-bundle, no link variables \(U_\ell\in G\), and
no continuum Lagrangian density for \(SU(3)\) Yang--Mills coupled to quarks are
part of the dynamical definition. Therefore Wilson loops and the QCD static
potential are undefined, and lattice or continuum confinement diagnostics do
not apply.
\end{proof}

\section{Contraction depth and the No-Go}

The a priori charge-correcting estimate for the completed map \(T^\sharp\)
continues to give a short mean contraction depth of order
\(N_\star=14\). The No-Go Theorem continues to hold: residue structure, tower
topology and flux democracy do not determine the long resonant structure. The
long trajectories of length \(539\) continue to require the external
fixed-count, holographic-window and phase-locking constraints already
identified by that theorem.

\section{No derivation of \(539\) or \(539.9\) from QCD}

\begin{theorem}[QCD does not supply \(539\) or \(539.9\)]
\label{thm:qcd-not-539}
Confinement in QCD is a well-established non-perturbative phenomenon of a
four-dimensional gauge theory. It does not supply a derivation of the integer
\(539\) or of the period \(539.9\) inside the present number-theoretic
construction.
\end{theorem}

\begin{proof}
Any derivation of \(539\) or \(539.9\) from QCD confinement would require a
map from the gauge-theoretic data (Wilson loops, string tension, mass gap) into
the ternary orbit length or phase period. No such map is given by residue
structure, tower topology, flux democracy, or \(T^\sharp\). Empirical and
lattice facts about QCD therefore remain external to the number-theoretic
invariants of the model.
\end{proof}

\begin{corollary}[Surrogates unchanged]
Optional surrogate loop weights or a surrogate string tension along constrained
ternary trajectories remain descriptive statistics of modular projections. They
neither implement QCD confinement nor lift the No-Go
(Chapter~\ref{ch:wilson-surrogate}).
\end{corollary}

% end
